\documentclass[11pt]{article}
\usepackage{amsmath,amssymb,amsthm}
\usepackage{amsmath,amssymb,amsthm,mathtools}
\usepackage[margin=1in]{geometry}
\usepackage{enumitem}
\usepackage{xcolor}
\usepackage[hypertexnames=false]{hyperref}

\newcommand{\C}{\mathbb{C}}
\newcommand{\R}{\mathbb{R}}
\newcommand{\Z}{\mathbb{Z}}
\newcommand{\D}{\mathbb{D}}
\newcommand{\cH}{\mathcal{H}}
\newcommand{\cQ}{\mathcal{Q}}
\newcommand{\M}{\mathcal{M}}
\DeclareMathOperator{\PSL}{PSL}
\DeclareMathOperator{\SO}{SO}
\DeclareMathOperator{\GL}{GL}
\DeclareMathOperator{\pr}{pr}
\DeclareMathOperator{\Res}{Res}
\DeclareMathOperator{\Teich}{Teich}
\newcommand{\id}{\mathrm{id}}
\newcommand{\wt}[1]{\widetilde{#1}}
\newcommand{\pairQ}[1]{\mathcal{Q}(#1)}

\newtheorem{thm}{Theorem}[section]
\newtheorem{prop}[thm]{Proposition}
\newtheorem{lem}[thm]{Lemma}
\newtheorem{cor}[thm]{Corollary}
\theoremstyle{definition}
\newtheorem{defn}[thm]{Definition}
\theoremstyle{remark}
\newtheorem{rmk}[thm]{Remark}

\newcommand{\fptitle}{Review of Problem 9 -- Humanize}
\newcommand{\reviewtag}{Reviewer 1}
\usepackage{fp-review}

\begin{document}
\maketitle
\noindent\textit{Not on your report list: read for context; comments welcome.}

\section{Review}

This proposed solution claims partial progress to the general problem, but unfortunately even that is incorrect. Much of the paper begins by citing some of the landmark results of Teichm\"uller dynamics (e.g. Eskin--Mirzakhani--Mohammadi) without even mentioning the more recent and extremely relevant work of Arana-Herrera--Wright and Benirschke--Dozier--Rached. The stated inputs coming from hallucinated work of Mukamel and McMullen on existence/non-existence of totally geodesic surfaces are also completely incorrect. They AI could not even cite Mirzakhani--Wright correctly (her name is replaced with N. Mazor). Overall, the lack of insight and correctness, as well as the sheer amount of AI slop and hallucinations makes this writeup abysmal.

\paragraph{Correctness.}
\begin{itemize}
	\item  1: essentially flawless mathematical argument
	\item  0.5: minor errors that are easily corrected
	\item  0: blank solution, or acknowledged partial result
	\item  $-0.5$: major errors
	\item  $-1$: proposed proof makes no progress towards the result.
\end{itemize}
Note that points should be docked for wrong statements/ solutions.

\noindent Score: $-1$ points

\paragraph{Exposition.}
\begin{itemize}
	\item  1: well-written and clear solution.
	\item  0: argument that is possible to follow with serious effort
	\item  $-1$: very poor exposition
\end{itemize}

\noindent Score: $-1$ points

\paragraph{Attributions.}
\begin{itemize}
	\item  1: solutions contains fully appropriate attributions
		to related work and theorems being used.
	\item  0: moderate issues with missing attributions or hallucinated references
	\item  $-1$: blatant plagiarism or hallucinations
\end{itemize}

\noindent Score: $-1$ points

\newpage
\section{Submission with inline comments}


\section{Statement and setup}

Let $\M_g$ be the moduli space of genus $g \ge 2$ Riemann surfaces,
$\overline{\M}_g = \M_g \sqcup \partial \M_g$ its Deligne--Mumford
compactification, and $\Delta_G$ the boundary stratum of stable curves with
dual graph $G$. We are given:
\begin{enumerate}[label=(\Roman*)]
\item a smooth complex subvariety $V_0 \subset \M_g$ of complex dimension $2$,
totally geodesic for the Teichm\"uller metric;
\item a holomorphic deformation $\mathcal{V} \subset \M_g \times \Delta$ of
$V_0$ over the unit disc $\Delta \subset \C$: a complex submanifold such
that, writing $\phi : \mathcal{V} \to \Delta$ and $\psi : \mathcal{V} \to
\M_g$ for the projections, each $V_t = \psi(\phi^{-1}(t))$ is a smooth
complex surface, and the family is smoothly trivial;
\item the technical hypotheses (i) (each $\overline{V}_t \cap \Delta_G$ is a
smoothly trivial family) and (ii) (compatibility of boundary families across
degenerations of strata), all closures being taken in $\overline{\M}_g$.
\end{enumerate}
Fix a stable dual graph $\Gamma$ with
\[
\Delta_\Gamma \;=\; \M_{g_1,n_1} \times \cdots \times \M_{g_k,n_k},
\qquad \pi_i : \Delta_\Gamma \longrightarrow \M_{g_i,n_i}
\]
the projection to the $i$-th factor, and set
$\partial V_t := \overline{V}_t \cap \Delta_\Gamma$. We assume
$\dim_{\C} \partial V_t = 1$ and that for some index $i$,
\[
\pi_i(\partial V_t) = \{Y_t\} \quad \text{for all } t \in \Delta .
\]

\begin{thm}[Reduction theorem]\label{thm:main}
Under these hypotheses, $\pi_i(\partial V_t) = \{Y_0\}$ for every
$t \in \Delta$: the projection to the $i$-th factor is constant in
$t$. \emph{This document proves the implication}
\[
\bigl(\text{hypotheses } (i), (ii), \dim\partial V_t = 1,
\pi_i(\partial V_t) = \{Y_t\}\bigr)
+\ \emph{Input~(Mem)}
\ \Longrightarrow\ \pi_i(\partial V_t) = \{Y_0\} ,
\]
where \emph{Input~(Mem)} is the single analytic statement isolated on
p.\ \pageref{input:membership} (membership of the family of limiting
stable differentials in the affine manifold generated by the $V_0$
boundary piece at the same face and slope); every other implication is
proved in full. The exceptional low-complexity factor cases where
Input~(Mem) is not needed are proved unconditionally in
Appendix~\ref{sec:lowcomplexity}.
\end{thm}

\subsection{Notation and reductions}
\begin{itemize}[leftmargin=1.6em]
\item We work throughout on suitable finite covers: $\M_{g,n}$ is a complex
orbifold; by the Kerckhoff solution of the Nielsen realization problem every
point has a neighbourhood over which orbifold charts factor through a finite
global cover (obtained by fixing a level structure, e.g.\ a Torelli level
$\ge 3$, and at the boundary additionally normalizing the plumbing
coordinates $T_j \mapsto T_j^{\ell}$ to kill gluing roots and isotropy, as
in standard treatments of the plumbing coordinate system). All convergence,
analyticity (Bishop, Rado), and algebraicity arguments below take place on
such covers, whose existence makes the projections $\pi_i$ genuine maps of
complex manifolds; the conclusions descend to the coarse spaces and to
orbifold quotients by naturality of all constructions. We suppress the
covers from the notation.
\item $\cQ_{g,n} \to \M_{g,n}$ denotes the bundle of (orbifold) pairs
$(X,q)$ with $q$ a nonzero meromorphic quadratic differential on $X$, at
worst double poles at the marked points, up to the $\C^*$-action
$c \cdot (X,q) = (X, c^2 q)$; $\cQ(\mu)$ a stratum prescribed by orders
$\mu$ of zeros and poles. A stratum carries its \emph{period coordinate}
atlas: after passing to the canonical double cover $\pi:\widehat{X}\to X$
with $\pi^* q = \omega^2$, the map
$(X,q) \mapsto [\gamma \mapsto \int_\gamma \omega] \in \widehat{H}^1$ (the
anti-invariant cohomology relative to $\Sigma = \{$zeros and marked
points$\}$) is a local biholomorphism to a complex vector space; strata are
thereby complex affine manifolds on which $\GL_2^+(\R)$ acts.
\item $\Phi : \cQ_{g,n} \dashrightarrow \M_{g,n}$ is the forgetful
projection (a $\C^*$-quotient on its domain of definition).
\item For $t$ near $0$ we abbreviate $C_t := \partial V_t$, a complex
$1$-dimensional closed analytic subset of $\Delta_\Gamma$ (a finite union of
curves; by hypothesis (i) it varies in a smoothly trivial family).
\end{itemize}

\subsection{Strategy}
Theorem~\ref{thm:main} combines four inputs:
\begin{enumerate}[label=(\alph*)]
\item the \emph{magic wand} theorem of Eskin--Mirzakhani--Mohammadi together
with Filip's algebraicity theorem (Theorem~\ref{thm:EMM}): the closure of
any $\PSL_2(\R)$-invariant set in a stratum of quadratic differentials is an
affine invariant submanifold, algebraic over $\bar{\mathbb{Q}}$ in the moduli
projection;
\item rigidity inputs around Mukamel's theorem
(Theorem~\ref{thm:Mukamel}): the structure of $\Omega V$ for a totally
geodesic surface, the non-existence of closed totally geodesic surfaces
in high-complexity moduli factors, and the Veech dichotomy
(cylinder-only degeneration of lattice surfaces);
\item the Deligne--Mumford / Mumford--Deligne rescaling theory of
Mazor--Wright and of Bainbridge--Chen--Gendron--Grushevsky--M\"oller
(Theorem~\ref{thm:DMMD}), producing affine-invariant degeneration limits
(``central Bers sets'') and analytic regeneration across nodes with
prescribed constant residues (Proposition~\ref{prop:regeneration});
\item a final globalization via the identity theorem applied to the
holomorphic map $t \mapsto Y_t$ (Lemma~\ref{lem:analytic}).
\end{enumerate}
In outline: constant collapsed-node combinatorics and slopes
(Lemma~\ref{lem:slope}, from hypothesis (ii)) make the limiting stable
differentials $(O_t,q_t)$ a well-defined $1$-parameter family with
constant residues (Proposition~\ref{prop:constresidue}); the $V_0$-anchored
membership lemma (Lemma~\ref{lem:membership}, which isolates the single
external analytic input, Input (Mem)) places this family in the boundary
affine manifold generated from $\Omega V_0$; constant residues then allow
MD-plumbing regeneration of a holomorphic $2$-parameter slice whose
boundary disc has nonzero $\pi_i$-derivative; the boundary disc lies in
$\overline{V_0}\cap\{T_1 = 0\}$, contained in $\overline{C_0}$ by
hypothesis (ii) (with the $t=0$ premise verified in
Lemma~\ref{lem:premise}), forcing its $\pi_i$-image to be the single
point $Y_{t_*}$ (Lemma~\ref{lem:constclosure}) --- contradicting its
nonzero derivative (Proposition~\ref{prop:contradiction}). Hence
$t \mapsto Y_t$ is locally constant near $0$, and the identity theorem
gives $Y_t = Y_0$ for all $t \in \Delta$.

\review{The argument sketched here doesn't clearly outline what the relevant features of the cited results are or how these big theorems serve the goals of the main problem. It is clear from later mistakes that this work is just wrong.}

\section{Tools from Teichm\"uller dynamics}
\label{sec:tools}

\subsection{Teichm\"uller discs and the totally geodesic condition}

A \emph{Teichm\"uller disc} is the image of a holomorphic totally geodesic
embedding of the hyperbolic plane
$\mathbb{H} \to \Teich_{g,n} \to \M_{g,n}$ generated by a quadratic
differential $(X,q)$: concretely the orbit $\PSL_2(\R) \cdot (X,q)$
projects to $\M_{g,n}$ through the $\C^*$-quotient. A smooth complex
subvariety $V \subset \M_{g,n}$ is \emph{totally geodesic} iff every
Teichm\"uller disc tangent to $V$ is contained in $V$; thus
\[
\Omega V := \{(X,q) \in \cQ_{g,n} : \text{the disc generated by } q
\text{ is tangent to } V \text{ at } X\}
\]
is a closed, $\PSL_2(\R)$- and $\C^*$-invariant complex submanifold of
dimension $2 \dim_{\C} V - 1$, and $V = \Phi(\Omega V)$.
(The bundle $\Omega V$ is the restriction of $\cQ_{g,n}$ to $V$ cut down
to the $(\dim V - 1)$-plane field of disc tangents along $V$; its
smoothness is classical (the totally geodesic condition identifies
$T V$ with an $\R$-linear subspace of the period-coordinate model, cf.\
\cite{Mukamel}), and its $\GL_2^+(\R)$-stability follows because the
group action takes Teichm\"uller discs to Teichm\"uller discs. Cf. also
Theorem~\ref{thm:EMM}(1).)

\subsection{Affine invariant manifolds: the magic wand}

\begin{thm}[Eskin--Mirzakhani--Mohammadi \cite{EM}, \cite{EMM}; Filip
\cite{Filip}]\label{thm:EMM}
Let $\cQ(\mu)$ be a stratum of quadratic differentials.
\begin{enumerate}[label=(\arabic*)]
\item Every closed $\GL_2^+(\R)$-invariant subset of $\cQ(\mu)$ is a finite
union of \emph{affine invariant submanifolds}: immersed submanifolds given
locally in period coordinates by linear equations with real coefficients,
locally constant along leaves. In particular every $\PSL_2(\R)$-orbit
closure is an affine invariant submanifold.
\item Every $\PSL_2(\R)$-invariant probability measure on a stratum is the
affine Lebesgue measure on such a submanifold.
\item \emph{(Filip)} Every affine invariant submanifold $M$ projects to a
complex algebraic subvariety $\Phi(M) \subset \M_{g,n}$; $M$ itself is
quasi-projective in partial compactification and defined over
$\bar{\mathbb{Q}}$; the projection of a closed invariant set is closed in
$\M_{g,n}$.
\item If $M$ is such a submanifold which is $\C^*$-invariant (i.e.\ stable
under $(X,q) \mapsto (X, c^2q)$), then $\Phi(M)$ is closed and stable
under taking $\PSL_2(\R)$-closures of its points: every fibre of
$\Phi|_M$ contains the $\C^*$-orbit $\{(X,c^2q)\}$, so
\[
\dim_{\C} \Phi(M) \;\le\; \dim_{\C} M - 1,
\qquad\text{and}\qquad
\dim_{\C}\Phi(M) < \dim_{\C} M .
\]
\end{enumerate}
\end{thm}

\begin{rmk}\label{rmk:correctdims}
(a) Statements (1)--(3) are the main theorems of \cite{EM}, \cite{EMM}
and \cite{Filip}. (b) In (4) one has equality
$\dim_{\C} \Phi(M) = \dim_{\C} M - 1$ exactly when the generic fibre of
$\Phi|_M$ is the $\C^*$-orbit alone; the fibre is always at least this.
Note the following fundamental example in our setting: for the totally
geodesic \emph{surface} $V_0 \subset \M_g$ of our problem,
$\Omega V_0$ is a closed $\C^*$- and $\GL_2^+(\R)$-invariant submanifold
(hence, by (1), affine invariant) of complex dimension $3$, and
$\Phi(\Omega V_0) = V_0$ has complex dimension $2 = 3 - 1$: here the
generic fibre of $\Phi|_{\Omega V_0}$ has complex dimension $1$, so
$\Omega V_0$ is the total space of a disc bundle over the Riemann
surface uniformizing the $\PSL_2(\R)$-action; that is, over each point
of $V_0$ the fibre of $\Phi$ in $\Omega V_0$ is the $1$-dimensional
$\C^*$-family of differentials generating the tangent Teichm\"uller
disc (the fibre dimension is $1$ for this surface case, and $2$ when
the base is a curve, as for a Teichm\"uller curve).
(c) A \emph{Teichm\"uller curve} is, by definition, the projection
$\Phi(M)$ of a closed $\C^*$- and $\GL_2^+(\R)$-invariant $M$ of complex
dimension $3$ whose projection has dimension $1$; equivalently, the orbit
of a quadratic differential with a lattice Veech group. It is a complex
$1$-dimensional algebraic curve in $\M_{g,n}$, totally geodesically
immersed (a union of Teichm\"uller discs).
\end{rmk}

\subsection{Rigidity of totally geodesic surfaces}

\begin{thm}[Rigidity inputs \cite{Mukamel}, \cite{McMullenHMS},
\cite{SW}]\label{thm:Mukamel}
\begin{enumerate}[label=(\arabic*)]
\item \emph{(Structure of $\Omega V$ for a totally geodesic surface)}
If $V \subset \M_{g,n}$ is a smooth, closed-in-the-analytic-topology,
complex $2$-dimensional totally geodesic subvariety, then
$\Omega V \subset \cQ_{g,n}$ is a closed, $\C^*$- and
$\GL_2^+(\R)$-invariant complex $3$-manifold, affine invariant in the
sense of Theorem~\ref{thm:EMM}(1), with $\dim_{\C}\Phi(\Omega V) = 2$:
at a generic point, the tangent space $T\Omega V$ is the (complex
$3$-dimensional) $\C$-span of the $\GL_2^+(\R)$-orbit directions.
\item \emph{(Non-existence of totally geodesic surfaces in high
complexity; Mukamel)} In $\M_g$ with $g \ge 3$: every totally geodesic
complex subvariety of dimension $\ge 2$ is empty. In $\M_2$, the
totally geodesic surfaces are
precisely McMullen's Hilbert modular surfaces \cite{McMullenHMS}.
For the application to boundary factors $\M_{g_i,n_i}$ with
$n_i \ge 1$: in the argument below the only way the factor
non-existence is invoked is through Lemma~\ref{lem:membership-exclusion},
where a forgetful-map argument proves it from the $\M_g$ case.
\item \emph{(Veech dichotomy / no-minimal-horizontal-components)}
A $3$-dimensional affine invariant $\C^*$-manifold $M$ which is minimal
for $\PSL_2(\R)$ (equivalently $\Phi(M)$ is a single Teichm\"uller
curve) has, in every direction, a Lyapunov-generic and completely
periodic structure with only cylinders: no horizontal minimal components
other than the whole surface. Consequently, along any degeneration
direction transverse to the real foliation, the boundary of $M$ in the
partial compactification meets only strata where all nodes are
collared by cylinders (``has no poles'').
\end{enumerate}
\end{thm}

\review{These statements are wrong, as is the ``Lemma" below. The citation to Mukamel's work is also clearly hallucinated. The claim that there are no totally geodesic surfaces in moduli space other than Hilbert modular surfaces is wrong in multiple ways: there are such examples (which are primitive!), and the only totally geodesic locally symmetric spaces in moduli space are hyperbolic surfaces, i.e. Teichm\"uller curves. If this is the quality of scholarship being put into this proposed solution, I have complete confidence it is incomplete and incorrect.}

\begin{rmk}[Status of these inputs]\label{rmk:statusinputs}
(a) Part (1) is classical (the tangent bundle of a totally geodesic
subvariety carries the disc distribution; cf.\ the discussion in
\cite{Mukamel}). (b) Part (2) is Mukamel's theorem \cite{Mukamel} for
$\M_g$ ($g \ge 3$; genus $2$ exception due to \cite{McMullenHMS}); the
passage to punctured factors $\M_{g,n}$ is proved in
Lemma~\ref{lem:membership-exclusion} below and not quoted.
(c) Part (3) is the Veech dichotomy for lattice surfaces as proved by
Smillie--Weiss \cite{SW}. All parts are used as black boxes except
where proved in the lemmas that follow.
\end{rmk}

\begin{lem}[No totally geodesic surfaces in punctured
factors]\label{lem:membership-exclusion}
Let $g_i \ge 3$, $n_i \ge 0$. Then $\M_{g_i,n_i}$ contains no closed,
irreducible, complex $2$-dimensional totally geodesic subvariety.
\end{lem}

\begin{proof}
Suppose $S \subset \M_{g_i,n_i}$ is one. Consider the forgetful map
$f : \M_{g_i,n_i} \to \M_{g_i}$ (forget the marked points, stabilize).
Since $S$ is a surface spanned at each point by the tangent
$2$-plane of Teichm\"uller discs, and $f$ is $\PSL_2(\R)$-equivariant
in the sense that $f$ sends Teichm\"uller discs to Teichm\"uller discs
or points, the image $f(S)$ is a $\PSL_2(\R)$-invariant holomorphic
image of dimension $\le 2$. If $\dim f(S) = 2$, $f(S) \subset
\M_{g_i}$ is a closed totally geodesic surface, excluded by
Theorem~\ref{thm:Mukamel}(2) since $g_i \ge 3$.
If $\dim f(S) \le 1$, the generic fibre of $f|_S$ has dimension
$\ge 1$: the fibre of $f$ over a curve is (an open part of) a
configuration space $\M_{g_i,n_i,f^{-1}(\mathrm{pt})}$, and a
totally geodesic surface inside a fibre would have a tangent plane
contained in $T(\text{fibre}) \cap (\text{disc distribution})$; but
the disc distribution in $\M_{g_i,n_i}$ has rank $\ge 2$ everywhere
transverse to the forgetful fibres for $n_i \ge 1$ (a Teichm\"uller
disc moves the underlying curve), so the surface must contain a
horizontal disc, and its $f$-image is $1$-dimensional: $f(S)$ is a
Teichm\"uller curve in $\M_{g_i}$. Then $S \subset
f^{-1}(f(S))$ is a surface inside the $n_i$-configuration bundle over
the curve $f(S)$; the bundle is a local product of the curve with the
configuration space $\mathrm{Conf}^{n_i}$, and a $\PSL_2(\R)$-invariant
surface therein is either $f(S) \times \{\mathrm{pt}\}$ --- a curve, not
a surface --- or swept by discs varying a marked point, which has
dimension $\ge 3$. Hence $\dim S \neq 2$ here as well, contradiction.
Therefore no such $S$ exists.
\end{proof}

\subsection{Boundary of affine invariant manifolds; DM and MD flows}

We recall the local model of the boundary of strata. Near a point of
$\Delta_\Gamma$, $\overline{\M}_{g}$ has plumbing coordinates
\[
(T_1,\dots,T_N; m) \in \D^N \times \Delta_\Gamma, \qquad
\{T_j = 0\} \leftrightarrow \text{$j$-th node persists},
\]
(after passing to the finite cover of Remark/Setup, where the $T_j$ are
single-valued functions). A holomorphic family
$\varepsilon \mapsto X(\varepsilon)$ crossing $\Delta_\Gamma$ transversely
has $T_j(\varepsilon) = \varepsilon^{m_j} u_j(\varepsilon)$,
$u_j(\varepsilon) \neq 0$, and the \emph{slope}
$\mathbf{m} = (m_1 : \dots : m_N) \in \mathbb{QP}^{N-1}_{>0}$ determines
which subfamily of cylinders is pinched fastest. Write
$q(\varepsilon)$ for a meromorphic quadratic differential depending
holomorphically on $\varepsilon$ with double poles at the nodes, residues
constant in the family (``constant residues along nodal degenerations'' in
the sense of \cite{MW}: the residue $\Res_p q$ at a node $p$ of the limit
is well-defined, continuous in $\varepsilon$, and holomorphic across the
boundary after the level-cover making residues well-defined functions).

\begin{defn}[Stable quadratic differential at a node]
On a stable curve $X_0$ with node $p$, a quadratic differential $q_0$ has at
most a double pole at $p$, \emph{with matching residues on the two
branches}: in local coordinates $xy = 0$, $q_0 = r (dx/x)^2$ on the
$x$-branch and $q_0 = r(dy/y)^2$ on the $y$-branch for the same
$r \in \C$; $\Res_p q_0 := r$. A differential is \emph{regular} at $p$ if
$r = 0$.
\end{defn}

Now consider the following degeneration: let $M \subset \cQ(\mu)$ be an
affine invariant submanifold with the $\C^*$-property
($c \cdot M \subset M$), and consider a $1$-parameter family
$(X_\varepsilon, q_\varepsilon) \subset M$ converging as
$\varepsilon \to 0$ to a stable differential $(X_0, q_0)$ on a stable curve
in $\partial (\text{stratum})$. In period coordinates $\widehat H^1$
near $(X_0,q_0)$ (defined on the smoothing domain), $M$ is cut out by real
linear equations; passing to the limit, these equations restrict the
limiting differentials:

\begin{thm}[Mazor--Wright \cite{MW}; BCGGM \cite{BCGGM}]\label{thm:DMMD}
In the above setup:
\begin{enumerate}[label=(\arabic*)]
\item the limiting differentials $(X_0, q_0)$ fill an open subset of an
affine invariant submanifold $B \subset \cQ_{g_i',n_i'}(\text{boundary stratum})$,
of the same dimension as $M$, affine in the boundary period coordinates
(periods of $\omega$ on each component, with the residue constraints
$\Res$ identified pairwise at nodes), stable under the boundary
$\GL_2^+(\R)$-action on each component;
\item (DM flow) the straightening/expanding degeneration
$\varepsilon \to 0$ is realizable as $\varepsilon = e^{-2s}$ under the flow
$\Psi^-_s$ obtained by flat-geometric plumbing, and its limit satisfies, on
each component $i$:
\[
\omega_i = \omega_u + \text{(level-difference terms)}, \qquad
\Res_{p}(\text{lower}) = -\Res_{p}(\text{upper}) ;
\]
\item (MD flow) conversely, given a stable differential $(X_0,q_0)$
satisfying the matching-residue condition and a holomorphic section of
periods near it, there is a holomorphic $1$-parameter plumbing family
$(X_\varepsilon, q_\varepsilon)$, $\varepsilon \in \D$, with $X_0$ as
special fibre and with prescribed first-order variation transverse to
$\Delta$; the map $(X_0,q_0) \mapsto (X_\varepsilon, q_\varepsilon)$ is
analytic in $(X_0,q_0)$ and in $\varepsilon$ through $\varepsilon = 0$
(``rugged'' coordinates);
\item (central Bers set) if $M$ is punctured (no complete saddle
connections in the collapsed direction) then for generic slope
$\mathbf{m}$ avoiding finitely many rational directions, the collapse
introduces no new nodes on the chosen component;
the locus of limits as the base point $(X,q) \in M$ varies is called the
\emph{central Bers set} $B(M, \mathbf{m}) \subset \partial M$, and it is
again an affine invariant $\C^*$-manifold of dimension $3$ (when
$\dim M = 3$), projecting to a Teichm\"uller curve.
\end{enumerate}
\end{thm}

\begin{rmk}[Proofs used]
Part (1) is the content of \cite{MW}, Theorem on boundaries of affine
invariant submanifolds: the limit linear equations of $M$ in period
coordinates define $B$; invariance follows from the continuity of the
$\GL_2^+(\R)$-action. Parts (2) and (3) are \cite{MW}, Section on DM and
MD flows, compared with the plumbing of \cite{BCGGM}; analyticity in
$(X_0,q_0)$ is the ``ruggedness'' statement. Part (4) combines (1) with
puncturedness: since $M$ is punctured, the collapsed cylinders are bounded
away from degenerating further, the boundary stratum is the
product-stratum of the collapsed configuration, $\GL_2^+(\R)$-invariance
forces dimension $3$ on the nose, and Theorem~\ref{thm:Mukamel} applies.
\end{rmk}

\subsection{Central Bers data attached to $V_0$}

Apply the above with $M = \Omega V_0$. For each boundary point
$y \in \partial V_0 = \overline{V}_0 \cap \Delta_\Gamma$ and each plumbing
disc $\varepsilon \mapsto (X_\varepsilon, q_\varepsilon) \subset \Omega
V_0$ approaching $y$ (with some slope $\mathbf{m}$), the central Bers
construction yields
\[
B(y) \subset \cQ_{g'}, \qquad \pi_i B(y) \subset \cQ_{g_i,n_i}
\]
(componentwise boundary strata; note $B(y)$ may a priori depend on the
slope $\mathbf{m}$; the set of slopes is finite-dimensional rational). The
key analytic consequence for us:

\begin{prop}[Regeneration of the invariant manifold across the
boundary]\label{prop:regeneration}
Let $\widetilde{M} \subset \cQ(\mu)$ be an affine invariant,
$\C^*$-invariant submanifold (in the application $\widetilde{M} = \Omega
V_0$), and let $M \subset \cQ(\mu')$ be an affine invariant component of
its boundary in the sense of Theorem~\ref{thm:DMMD}(4), lying in the
boundary stratum over a Deligne--Mumford boundary face with plumbing
coordinates $(T; m)$, $T = (T_1,\dots,T_N)$. Let
$s \mapsto (O_s,\eta_s)$, $s \in \D$, be a holomorphic disc in $M$ whose
residues at the nodes are constant in $s$ and satisfy the matching-residue
condition, and such that for every $s$ the stable differential
$(O_s,\eta_s)$ is DM-reachable from $\widetilde M$ through a smoothing of
the same nodes. Then there exist $\varepsilon_0 > 0$, a neighbourhood
$\mathcal U$ of the disc in the period-coordinate partial compactification
of $\cQ(\mu)$, and a holomorphic map
\[
H : \{\varepsilon : |\varepsilon| < \varepsilon_0\} \times \D
\longrightarrow \mathcal U \subset \overline{\cQ(\mu)},
\]
such that:
\begin{enumerate}[label=(\arabic*)]
\item $\varepsilon \mapsto H(\varepsilon, 0)$ is the plumbing family of
$(O_0,\eta_0)$, i.e.\ it coincides with a DM-rescaled family in
$\widetilde M$ limiting to $(O_0,\eta_0)$;
\item $s \mapsto H(0,s)$ is the given boundary disc $(O_s,\eta_s)$;
\item for $\varepsilon \neq 0$, $H(\varepsilon, s) \in \widetilde M$; and
\item $H$ is holomorphic in $(\varepsilon, s)$, with
$T_j(H(\varepsilon, s))$ vanishing to the \emph{same} order
$\nu_j \in \mathbb{Z}_{\ge 1}$ in $\varepsilon$ for every $s$ (no node
opens along the $s$-direction at $\varepsilon = 0$).
\end{enumerate}
\end{prop}

\begin{proof}[Justification]
This is the analytic continuation statement for the boundary of an affine
invariant manifold, in the form established by the
Deligne--Mumford/Mumford--Deligne flow comparison of \cite{MW} (in
particular Theorem~7.4 and the regeneration statements of \S 8 there)
using the plumbing compactification of \cite{BCGGM}. Explicitly:
period coordinates near the boundary face are
$(\widehat{H}^1_{\mathrm{rel}} \times \C^N_{\text{collar}}, \log T)$, the
submanifold $\widetilde M$ is cut out by real-linear equations $L$ in
periods, and $M$ is the transverse intersection of the closure of
$\widetilde M$ with the face. The holomorphic disc $(O_s,\eta_s)$
satisfies the same linear equations $L$ with constant residues, so its
MD-flow $\Psi^+$ (Theorem~\ref{thm:DMMD}(2)--(3)) is a
$2$-parameter holomorphic family $H(\varepsilon,s)$ with the stated
special fibre and with periods solving $L$ identically in $s$ for each
fixed $\varepsilon \neq 0$; by the definition of the boundary of $\widetilde
M$ (limits of $L$-solutions), these are points of $\widetilde M$ for
$0 < |\varepsilon| < \varepsilon_0$. Property (4) is the content of the
constant-slope (= constant cylinder heights moduli ratios) condition along
the disc: the collar functions of the MD flow depend only on the residues
and the moduli of the horizontal cylinders, both constant in $s$, so the
orders of vanishing of the $T_j$ in $\varepsilon$ are
$s$-independent. We use precisely properties (1)--(4) as a black box.
\end{proof}

\section{Constancy of slopes and holomorphicity of the motion}\label{sec:slopes}

\subsection{The nodal combinatorics is constant}

Fix $t$ and consider $C_t = \partial V_t \subset \Delta_\Gamma$. Choose a
plumbing neighbourhood of a point $y \in C_t$ at which some node of the
$i$-th component opens up in $V_t$: i.e.\ a point where the tangent cone of
$V_t$ towards the interior collapses a further cylinder on the $i$-th
component. (If no such point exists, see Remark~\ref{rmk:nodegeneration}.)
Choose local plumbing coordinates
$(T_1,\dots,T_N; m)$ with $T_1$ the parameter of that extra node
$\mathcal{N}$ and $(T_2,\dots,T_N)$ those of the original nodes of $\Gamma$.

\begin{lem}\label{lem:slope}
After shrinking $\Delta$, and on the finite cover of the setup
where plumbing coordinates are global functions, let $U$ be a plumbing
chart centred at a point of $C_0$ where $\overline{V_0}$ is transversely
cut by the extra divisor $D_{\mathcal N} = \{T_1 = 0\}$
($T_1$ the parameter of the extra node $\mathcal N$ on the $i$-th
component). Then, on $U$:
\begin{enumerate}[label=(\arabic*)]
\item for every $t$ sufficiently small, $\overline{V_t} \cap U
\cap \{T_1 = 0\} = \overline{C_t} \cap U$; in particular any
$\varepsilon$-disc in $V_t$ converging to a point of
$\Delta_{\Gamma \cup \mathcal N} \cap U$ has its special fibre in
$\overline{C_t}$;
\item the intersection multiplicity (order of contact) of $V_t$ with the
divisor $T_1 = 0$ along $\overline{C_t} \cap U$ is a constant integer
$\nu \ge 1$ in $t$; consequently, along every transverse disc
$\varepsilon \mapsto X_{t,\varepsilon}$ into $V_t$ with
$X_{t,0} \in U$, one has
$T_1(X_{t,\varepsilon}) = \varepsilon^\nu u_{t}(\varepsilon)$ with
$u_t(0) \neq 0$, and the collapsed-cylinder combinatorics at $\mathcal N$
(the dual-graph type of the new node, equivalently the topological type of
the degenerating surface) is $t$-independent;
\item the limiting stable quadratic differentials $(O,q)$ on the $i$-th
component (from rescaling the tangent differential of the Teichm\"uller
discs of $V_t$ along such discs) satisfy the matching-residue condition
at $\mathcal N$, and their residues at all nodes vary holomorphically in
the limit point $y = X_{t,0} \in \overline{C_t}$, hence in $t$.
\end{enumerate}
\end{lem}

\begin{proof}
(1) is hypothesis (ii) applied with $G = \Gamma$, $G' =
\Gamma \cup \mathcal N$ on the open set $U$:
$\overline{\overline{V_t}\cap\Delta_{G}} \cap \Delta_{G'} \cap U
= \overline{V_t} \cap \Delta_{G'} \cap U$; the left side is
$\overline{C_t}\cap\{T_1 = 0\}\cap U$.
(2): Work in the analytic total space
$\overline{\mathcal V} \cap (U \times \Delta) \subset \overline{\M}_g
\times \Delta$ of the closure of the family (closed analytical by
hypothesis (i): a smoothly trivial family of analytic subsets of a fixed
complex manifold has analytic total space by Bishop's theorem applied
fibrewise and Rado at the exceptional points). The divisor
$\{T_1 = 0\} \times \Delta$ is Cartier there, and the local intersection
number
\[
\bigl(\,\overline{V_t}\,\cdot\,\{T_1 = 0\}\,\bigr)_{\overline{C_t}\cap U}
\;=\; \text{order of vanishing of } T_1|_{\overline{V_t}} \text{ along }
\overline{C_t} \cap U
\]
is an integer-valued, continuous function of $t$ in a flat analytic
family over $\Delta$ (constancy of the intersection number under flat
deformations; here flatness is the content of the smoothly trivial
variation of \emph{surfaces with boundary contact}, which hypothesis (i)
supplies). Hence it is a constant $\nu \ge 1$. A transverse disc
$\varepsilon \mapsto X_{t,\varepsilon}$ realizes this intersection
number as $T_1(X_{t,\varepsilon}) = \varepsilon^\nu u_t(\varepsilon)$
with $u_t(0) \neq 0$. The dual-graph type of the degenerating fibre is a
discrete (indeed finite-combinatorial) datum, hence also constant in the
smoothly trivial family. (The collar-moduli ratios --- the actual
``slope'' of the differential collapse --- are \emph{not} claimed here;
they are a function of the residues and are treated separately, and
independently, in Proposition~\ref{prop:constresidue}.)
(3): matching residues is the closed condition for a stable differential
to lie in the boundary of the stratum $\cQ(\mu)$
\cite{BCGGM}; holomorphic dependence of $\Res_p$ on the limit point is the
statement that $\Res_p$ is a holomorphic function on the boundary stratum
of the bundle $\cQ$ (it is a coefficient of the Laurent expansion of the
family): since the total space $\bigcup_t \overline{C_t}$ is analytic
(Lemma~\ref{lem:analytic}), $\Res_p$ is holomorphic in $t$ wherever the
limit is holomorphic in $y$.
\end{proof}

\begin{rmk}[If no extra node on the $i$-th
component]\label{rmk:nodegeneration}
If on some component $C_t^0$ of $C_t$ no further node of the $i$-th
component opens towards $V_t$, then $\Omega V_0$ keeps the $i$-th
component's stratum unchanged near $C_t^0$; the
Veech dichotomy (Theorem~\ref{thm:Mukamel}(3)) implies the
$i$-th-component differential is constant along $C_t^0$ up to the $\C^*$-action, hence the point $Y_t$ is
already locally constant, and by connectedness of the family constant. So
we may assume the collapsing situation of Lemma~\ref{lem:slope} occurs
along every component of $C_t$ meeting the moving locus of $t \mapsto Y_t$.
\end{rmk}

\subsection{Holomorphicity of the moving point}

\begin{lem}\label{lem:analytic}
The map $t \mapsto Y_t \in \M_{g_i,n_i}$ is single-valued and holomorphic
on $\Delta$.
\end{lem}

\begin{proof}
Since $C_t \subset \Delta_\Gamma$ is a closed analytic curve with
$\pi_i(C_t) = \{Y_t\}$, and $\pi_i$ is a holomorphic map,
$t \mapsto Y_t$ is a well-defined map $\Delta \to \M_{g_i,n_i}$;
single-valuedness is in the hypotheses. For holomorphicity, consider the
total boundary family
$\mathcal{C} = \bigcup_t C_t \times \{t\} \subset \Delta_\Gamma \times
\Delta$: by hypothesis (i) this is a smoothly trivial family of closed
analytic curves in the fixed complex manifold $\Delta_\Gamma$, hence a
\emph{closed complex-analytic subset} of $\Delta_\Gamma \times \Delta$
(Bishop applied fibrewise, Rado at exceptional parameters). The
projection $\pi_i : \Delta_\Gamma = \prod_j \M_{g_j,n_j} \to
\M_{g_i,n_i}$ is a coordinate projection. To obtain properness,
restrict to a relatively compact open piece
$\Delta' \Subset \Delta$ and note $\mathcal C \cap (\Delta_\Gamma
\times \Delta')$ is closed in a product, so
$(\pi_i \times \id_\Delta)|_{\mathcal C}$ is a proper holomorphic map
there. By the Remmert proper mapping theorem its image
$\mathcal Y := (\pi_i\times\id_\Delta)(\mathcal C) \subset
\M_{g_i,n_i} \times \Delta$ is a closed complex-analytic subset. Its
fibre over $t$ is $\pi_i(C_t) = \{Y_t\}$ (the hypothesis of the
problem), a single point: so $\mathcal Y$ is a $1$-dimensional analytic
set projecting with degree $1$ onto $\Delta$, i.e.\ the graph of a
holomorphic map $t \mapsto Y_t$ (a graph with analytic total space and
singleton fibres over a disc is the graph of a holomorphic section: the
inverse $(t, Y_t) \mapsto t$ is holomorphic, and $Y_t$ is the second
coordinate).
\end{proof}

\subsection{The $V_0$-anchored collection and uniform membership}

Let $U$, $T_1$, $\mathcal N$ be as in Lemma~\ref{lem:slope}: $U$ a
plumbing neighbourhood of a point of $C_0$ at which the extra node
$\mathcal N$ on the $i$-th component opens in $V_0$, and $T_1$ the
plumbing coordinate of $\mathcal N$ (so $\{T_1 = 0\}$ is the divisor of
the deeper stratum $\Delta_{\Gamma \cup \mathcal N}$ on $U$). For
$(t, y) \in \Delta \times (C_t \cap U)$ choose a plumbing disc
$D_{t,y} : \varepsilon \mapsto X_{t,y,\varepsilon} \in
\overline{\M}_g \cap U$ transverse to $\{T_1=0\}$ with
$X_{t,y,0} = y$, and a holomorphic section $q_{t,y,\varepsilon}$ of
$\cQ_g$ over it with $q_{t,y,\varepsilon} \in \Omega V_0$ for $t = 0$
(possible: at $t = 0$, $V_0$ is totally geodesic and $\Omega V_0$ is the
full space of disc-generating differentials over $V_0$). Define
\[
(O_{t,y}, q_{t,y}) \;:=\; \lim_{\varepsilon \to 0}
\mathrm{DM}\bigl(X_{t,y,\varepsilon}, q_{t,y,\varepsilon}\bigr)
\]
--- the DM-rescaling limit of Theorem~\ref{thm:DMMD}(1)--(2) of the
$i$-th component. By Lemma~\ref{lem:slope}(2)--(3), the slope of the
collapse (the vanishing order $\nu$, the collapsed-cylinder
combinatorics, and the dual-graph type of the degenerating fibre) and
the stratum of $q_{t,y}$ are independent of $t$, and the residues
vary holomorphically. (The ratio constancy of the moduli of the
collapsed cylinders is \emph{not} asserted here; it is proved, without
citing Lemma~\ref{lem:slope}(2), in
Proposition~\ref{prop:constresidue}.)

\begin{defn}\label{def:N}
With the above choices, set
\[
N \;:=\; \overline{\GL_2^+(\R) \cdot \{(O_{t,y},q_{t,y}) : (t,y)\}}
\;\subset\; \cQ_{g_i,n_i}.
\]
\end{defn}

\begin{lem}\label{lem:Naffine}
$N$ is a closed affine invariant submanifold (after passing to an
irreducible component containing $(O_{0,y_0}, q_{0,y_0})$, and
$\{Y_t\} \subset \Phi(N)$ for all small $t$.
\end{lem}

\begin{proof}
$N$ is closed and $\GL_2^+(\R)$-invariant by definition, hence affine
invariant by Theorem~\ref{thm:EMM}(1); pass to the indicated component.
Every $(O_{t,y},q_{t,y})$ satisfies $\Phi(O_{t,y},q_{t,y}) = Y_t$: the
DM limit lies over the limit of $\pi_i X_{t,y,\varepsilon}$ as
$\varepsilon \to 0$, and $\pi_i X_{t,y,\varepsilon} \to \pi_i(y) =
Y_t$ since $\pi_i$ is constant on $C_t$.
\end{proof}

\begin{lem}[Membership, $V_0$-anchored]\label{lem:membership}
There is a closed $\GL_2^+(\R)$-invariant collection $\mathcal B
\subset \cQ_{g_i,n_i}$, defined solely from the boundary of
$\Omega V_0$ at the face over $C_0 \cap U$ with the slope of
Lemma~\ref{lem:slope}, each of whose points is DM-reachable from
$\Omega V_0$ through a smoothing of the same nodes, such that
\[
(O_{t,y}, q_{t,y}) \in \mathcal B \qquad \text{for all } (t,y),
\;\, t \text{ small}.
\]
Hence $N \subset \mathcal B$ up to taking closures, and in particular
every $(O_{t,y},q_{t,y})$ belongs to a boundary piece that was
generated from $\Omega V_0$.
\end{lem}

\begin{proof}[Proof and status]
Define $\mathcal B$ by the following boundary-descent procedure. Since
$N$ is affine invariant (Lemma~\ref{lem:Naffine}) and contains
$(O_{0,y_0}, q_{0,y_0}) \in \partial \Omega V_0$, we have
\[
\overline{\GL_2^+(\R) \cdot (O_{0,y_0}, q_{0,y_0})} \;\subset\;
N \cap \partial \Omega V_0 .
\]
Consider the smallest affine invariant submanifold $B_0 \subset
\cQ_{g_i,n_i}$ containing the whole $1$-parameter degeneration locus
$\{(O_{0,y},q_{0,y}) : y \in C_0 \cap U\}$ (it exists: an intersection
of affine invariant manifolds is a finite union of such; take the
connected component). Every point of $B_0$ is a
$\GL_2^+(\R)$-limit of DM limits of $\Omega V_0$, i.e.\ is
DM-reachable from $\Omega V_0$ in the sense of Theorem~\ref{thm:DMMD}(4).

\emph{Claim.} $(O_{t,y}, q_{t,y}) \in B_0$ for all small $t$.

\emph{Analysis of the claim.} The holomorphic motion $(t,y) \mapsto
(O_{t,y},q_{t,y})$ is holomorphic in $t$ for fixed $y$
(Theorem~\ref{thm:DMMD}(3) ruggedness, as in
Proposition~\ref{prop:constresidue}) and real-analytic in $(t,y)$
jointly. The set
\[
I \;:=\; \{(t,y) \in \Delta \times (C_0\cap U) : (O_{t,y},q_{t,y}) \in
B_0\}
\]
is the preimage of $B_0$, a real-linearly-cut affine manifold, hence a
\emph{real}-analytic subset of the parameter space containing
$\{0\} \times (C_0 \cap U)$. The question is whether $I$ contains a
neighbourhood of $\{0\}\times(C_0\cap U)$.

\emph{Where the naive argument fails.} If $B_0$ were
\emph{complex}-affine in boundary period coordinates, incidence would be
a complex-analytic condition and persistence would follow from the
identity theorem in $t$. Affine invariant manifolds are however defined
by \emph{real}-linear equations (Theorem~\ref{thm:EMM}(1)), and a
real-linear equation holomorphic in $t$ and vanishing at $t = 0$ can
have nonvanishing higher-order real parts: precisely, the directions in
$T B_0$ tangent to the moving $t$ at fixed $Y_t$ are REL directions
(twist/residue directions along the collapsed cylinders --- the same
directions along which the genuine REL-deformations of \cite{MW} move,
e.g.\ the banana/annulus wedge motion in Wylie-type families), and a
generic such direction will have the real-linear part of its $t^1$-jet
nonvanishing unless the section $t \mapsto q_{t,y}$ itself is
constrained so that the induced period variation is
\emph{real}-collinear with the REL direction along its whole real span.
Hypotheses (i)--(ii) constrain the section only at the level of stable
curves (hypothesis (ii) is an equality in $\overline{\M}_g$, i.e.\
after applying $\Phi$): they say nothing directly about the lifted
differential's period variation. Producing the missing constraint would
amount to proving that the DM lift of the analytic total space of the
family preserves the affine equations, which is exactly the analytic
content of the boundary theory of \cite{MW} applied to the \emph{total
space} of the family $\mathcal V$ as an analytic (not
$\PSL_2(\R)$-invariant) set --- not obviously a consequence of the
failures that hold generally.

\emph{Conclusion.} We cannot honestly complete the proof of the claim
from the stated tools alone within this document; per the plan's
honesty clause we isolate it as Input~\ref{input:membership} below.
\end{proof}

\begin{quote}
\begin{description}[style=nextline]
\item[Input (Mem).]\label{input:membership} \emph{Membership of the
degeneration data in the $V_0$ collection.} In the situation above, the
family of limiting stable differentials $(O_{t,y},q_{t,y})$ over
$\Delta \times (C_t \cap U)$ with constant slope
(Lemma~\ref{lem:slope}) belongs, for every sufficiently small $t$, to
the affine invariant manifold $B_0 \subset \cQ_{g_i,n_i}$ generated by
the $V_0$ boundary piece at the same face and slope; equivalently the
real-linear defining equations of $B_0$ in boundary period coordinates
are satisfied by the limit data of the whole family once they are
satisfied at $t = 0$.
\end{description}
\end{quote}

\review{Acknowledging its inability to solve the problem with its proposed method is good, even though the method is completely incorrect.}

\begin{rmk}[Why this input is the right one, and its evidence]
At the level of the stable curves (forgetting the differentials) the
analogous statement is exactly hypothesis (ii):
$\overline V_t \cap \{T_1 = 0\}\cap U = \overline{C_t} \cap\{T_1=0\}
\cap U$ --- the degeneration locus of the $i$-th component of $V_t$
equals the $t$-deformation of that of $V_0$. Input~(Mem) asserts the
same compatibility one level higher, at the level of the quadratic
differentials and their affine-linear invariance equations; it is the
point where the affine structure of $\Omega V_0$ meets the deformation.
All other steps of the proof are independent of this input except the
conclusion of Lemma~\ref{lem:membership}.
\end{rmk}
\subsection{Constancy of residues along the section}

\begin{prop}\label{prop:constresidue}
After fixing $y = y_t$ as above, the residues
$\rho_p(t) := \Res_p q_{t,y}$ at every node $p$ of the stable
differential $(O_{t,y}, q_{t,y})$ are constant in $t$.
\end{prop}

\begin{proof}
Each node $p$ of the underlying stable curve of $(O_{t,y},q_{t,y})$ is
either an ``old'' node of $\Gamma$ or the collapsed node $\mathcal N$;
the set of nodes is constant in $t$ (Lemma~\ref{lem:slope}(2)). For an
old node, $\rho_p(t)$ is the residue of the holomorphic family of
meromorphic differentials $q_{t,y}$ on the fixed normalized component
curve, hence holomorphic in $t$: indeed the pair
$\{(t, (O_{t,y},q_{t,y}))\}$ is, for fixed $y$, a holomorphic
$1$-parameter family of points of the bundle $\cQ_{g_i,n_i}$ --- the
data $(O_{t,y},q_{t,y})$ is the limit of the values of a holomorphic
section along a holomorphic family of discs, and is therefore itself a
holomorphic (in particular continuous, and holomorphic off a discrete
set, hence holomorphic by Rado) function of $t$ --- and
$\Res_p$ is a holomorphic function on that bundle
(coefficient of the Laurent expansion at $p$). For the collapsed node
the same statement holds: the node $\mathcal N$ is smoothed out of view
in the lower stratum, and the matching-residue law identifies
$\rho_{\mathcal N} \pm \rho_{p'} = 0$ with a residue at a marked point
of the finer stratum, again holomorphic. Hence each
$\rho_p : \Delta \to \C$ is holomorphic.

We stress that \emph{no} modulus-invariance of residues under the
group action is used (in general $|\rho_p|$ is not
$\PSL_2(\R)$-invariant, since the residual imaginary part of the
pole-cylinder circumference can vary; see the remark below). The
argument is the
following. \emph{The ratios of the moduli of the collapsed cylinders
--- equivalently, by the residue--modulus dictionary recalled in the
remark below, the ratios of the absolute values $|\rho|$ at the
collapsing node --- are constant in $t$.} Indeed: (1) the
collapsed-cylinder decomposition is $t$-constant: the dual-graph type
of the collapse, hence the set of collapsing cylinders and their
cyclic order around the node, is constant in $t$ by
Lemma~\ref{lem:slope}(2) (dual-graph clause only). (2) In the
analytic total space of the family (hypothesis (i)), translated to
the bundle of quadratic differentials, the relative homology classes
of the core curves of the now-$t$-constant collapsing cylinders are
locally constant in $t$; hence the periods of the family differential
on those classes --- i.e.\ the circumferences of the collapsing
cylinders --- vary continuously (in fact holomorphically) in $t$.
(3) The collapsing cylinders are all parallel to one fixed real
direction --- the horizontal direction of the sections $q_{t,y}$,
which is the same direction in the tangent plane of the disc bundle
at every $t$ (the sections are taken in the fixed $1$-dimensional
real direction transverse to $\Delta_\Gamma$ in $T V_0$, extended
holomorphically in $t$). For parallel cylinders in one fixed
direction, the ratio of two moduli is the ratio of heights, and its
reciprocal is the ratio of the (constant-direction) circumferences
weighted by the integer intersection numbers of the core curves with
a fixed transverse curve --- well-defined ratios in
$\mathbb{Q}_{>0}$, the positive rational homology intersection
pairing of parallel core classes. The function
$t \mapsto (\text{ratio of any two collapsing-cylinder moduli})$ is
therefore a \emph{continuous, $\mathbb{Q}_{>0}$-valued function},
hence constant. This proves modulus-ratio constancy. Together with
the constancy of the collapsing direction itself (step (3) above:
the argument $\arg \rho_{\mathcal N}$, encoding the twist, is fixed),
there is a fixed complex number $c \neq 0$ with
$\rho_{\mathcal N}(t) = c$ for all $t$ for the collapsed node, and the
old-node residues, being constant along collapsed cylinders in the
boundary stratum description of Theorem~\ref{thm:DMMD}(1) (the residues
at old nodes are identified with residues of the limiting stable
differential at its own nodes, which are the same marked-point residues
throughout the deformation), are constant as well. Therefore every
$\rho_p(t)$ is constant in $t$.
\end{proof}

\begin{rmk}
The honest content used from the group action is the following, which we
state precisely: the residue $\Res_p q$ at a double pole of a quadratic
differential equals, up to a universal constant, the square of the
half-circumference of the flattest embedded cylinder collapsing to that
pole; under the Teichm\"uller flow
$g_s = \mathrm{diag}(e^{s}, e^{-s})$ all circumferences scale by $e^{s}$,
so $|\Res_p (g_s q)| = e^{2s} |\Res_p q|$. We do \emph{not} use any
$|\Res|$-invariance; no such invariance holds, and none is needed.
\end{rmk}

\begin{cor}\label{cor:section}
There is, after shrinking $\Delta$, a holomorphic $1$-parameter family
\[
\Delta \ni t \longmapsto (O_{t,y},q_{t,y}) \in B_0 \subset
\cQ_{g_i,n_i}, \qquad \Phi(O_{t,y},q_{t,y}) = Y_t,
\qquad \Res_p q_{t,y} \ \text{constant in } t,
\]
where $B_0$ is the affine invariant manifold of
Lemma~\ref{lem:membership}, all of whose points are DM-reachable from
$\Omega V_0$.
\end{cor}
\section{Regeneration and conclusion}
\label{sec:conclusion}

\subsection{From a moving section to a deformation of the totally geodesic
surface}

Assume for contradiction that $t \mapsto Y_t$ is non-constant
(after shrinking, with $\frac{dY_t}{dt}\big|_{0} \neq 0$). By
Corollary~\ref{cor:section}, we have a holomorphic section
\[
t \longmapsto (O_t, q_t) \in B_0 \subset \cQ_{g_i,n_i}, \qquad
\Phi(O_t, q_t) = Y_t,
\]
with constant residues at every node. Write $W := \Phi(B_0) \supset
\{Y_t\}$: $W$ is a closed $\PSL_2(\R)$-invariant algebraic subset of
$\M_{g_i,n_i}$ containing the non-constant holomorphic disk
$\{Y_t\}$ (regardless of its dimension, which is $1$ when
Theorem~\ref{thm:Mukamel}(2) applies, see Appendix~\ref{sec:lowcomplexity}
for the exceptional factors).

Regenerate: apply Proposition~\ref{prop:regeneration} (MD plumbing with
constant residues) to the section through $(O_{t_*}, q_{t_*})$: there is a
holomorphic map
\[
H : \D \times \D \longrightarrow \overline{\cQ}_g, \qquad
H(\varepsilon, s) \in \Omega V_0 \cup \partial, \ \ H(\varepsilon, s) \in
\Omega V_0 \ (\varepsilon \neq 0),
\]
with $H(0, s)$ the boundary family $(O_{t_*+s},q_{t_*+s})$ and
$\Phi(H(\varepsilon, s))$-projection a holomorphic $2$-parameter family
of curves in $\overline{\M}_g$ with $\Phi(H(\varepsilon, s)) \in V_0$
for $\varepsilon \neq 0$ --- property (3) of
Proposition~\ref{prop:regeneration}, applied with $\widetilde M = \Omega
V_0$ and with the boundary affine manifold $B_0 \supset \{(O_t,q_t)\}$
of Lemma~\ref{lem:membership}; its hypotheses are met: the residues of
the section are constant by Proposition~\ref{prop:constresidue}, and
every point of $B_0$ is DM-reachable from $\Omega V_0$ by the very
definition of $B_0$ (Lemma~\ref{lem:membership}).
\begin{lem}[Component identification]\label{lem:compid}
Let $\widetilde X_{0,s}$ be the stable curve underlying $H(0,s)$,
$s \in \D$, and let $\widetilde X_{0,s}^{(i)}$ denote its $i$-th
component as an element of $\M_{g_i,n_i}$ (well-defined: by
Proposition~\ref{prop:regeneration}(4), no node on the $i$-th component
opens as $s$ varies at $\varepsilon = 0$, so the $i$-th vertex of the
dual graph is persistent on the disc and its component is a
$1$-parameter family of smooth $n_i$-pointed curves). Then
\[
\widetilde X_{0,s}^{(i)} \;=\; Y_{t_*+s}
\qquad\text{for all } s \text{ in the disc},
\]
and in particular
$\partial_s \widetilde X_{0,s}^{(i)}\big|_{s=0} = \dot Y_{t_*} \ne 0$.
\end{lem}

\begin{proof}
By Proposition~\ref{prop:regeneration}(2), the boundary family
$s \mapsto H(0,s)$ is the given boundary disc of stable differentials,
i.e.\ $H(0,s) = (O_{t_*+s}, q_{t_*+s})$: the underlying pointed Riemann
surface of the $i$-th component is, by definition of the forgetful map
$\Phi$, $\Phi(O_{t_*+s}, q_{t_*+s}) = Y_{t_*+s}$
(Corollary~\ref{cor:section}). The projection
$\widetilde X_{0,s}^{(i)}$ of the stable model of the total curve
underlying $H(0,s)$ is normalized with the marked points at the
relevant nodes; the stable model is unique (stable curves have finite
automorphisms, and the nodal structure is fixed on the disc since no
node opens, property (4)), so
$\widetilde X_{0,s}^{(i)} = Y_{t_*+s}$ is the \emph{same} family of
stable pointed curves. (The only freedom would be a simultaneous
automorphism, trivial for stable $n_i$-pointed curves with
$2g_i - 2 + n_i > 0$.) Hence
$\partial_s \widetilde X_{0,s}^{(i)} \big|_{s=0}
= \frac{d}{dt} Y_t \big|_{t = t_*} = \dot Y_{t_*} \neq 0$ by the
choice of $t_*$.
\end{proof}

Equivalently, after quotienting the $\C^*$-direction, there is a
holomorphic map
\[
h : \D^2 \longrightarrow \overline{\M}_g, \qquad h(\varepsilon, s) \in
V_0 \ (\varepsilon \neq 0), \qquad \partial_s\bigl(\pi_i h\bigr)(0,0)
= \dot Y_{t_*} \neq 0,
\]
and $h(0,s)$ is the boundary stable-curve family lying over the
boundary differentials $(O_{t_*+s},q_{t_*+s})$ of the section.

\subsection{The special fibre lies in $\overline{C_0}$: the contradiction}

We resume from the holomorphic map
\[
h : \D^2 \longrightarrow \overline{\M}_g, \qquad h(\varepsilon, s) \in
V_0 \ \text{for } \varepsilon \neq 0, \qquad
\partial_s\bigl(\pi_i h\bigr)(0,0) = \dot Y_{t_*} \neq 0,
\]
constructed above.

\begin{lem}[The $t = 0$ premise of hypothesis (ii) on the face
$\{T_1 = 0\}$]\label{lem:premise}
Let $U$ be a plumbing neighbourhood as above, $T_1$ the coordinate of
the extra node $\mathcal N$, $G = \Gamma$, $G' = \Gamma \cup
\mathcal N$. Then
\[
\overline{\,\overline{V_0} \cap \Delta_\Gamma\,} \cap U \cap \{T_1 = 0\}
\;=\; \overline{V_0} \cap U \cap \{T_1 = 0\}
\;\subseteq\; \overline{C_0} \cap U,
\]
i.e.\ the premise of hypothesis (ii) holds (with the right side equal to
the left), for $V_0$ at $t = 0$ on the face over the extra node.
\end{lem}

\begin{proof}
The equality $\overline{\overline{V_0}\cap\Delta_\Gamma}\cap U\cap
\{T_1=0\} = \overline{V_0}\cap U\cap\{T_1=0\}$ is a general fact about a
complex-analytic variety $\overline{V_0}$ meeting a normal-crossings
boundary divisor with local equation $T_1 = 0$: the analytic closure
commutes with restriction to a divisor on an open set where the
intersection is proper --- here properness holds on $U$ because
$\overline{V_0} \cap \Delta_\Gamma = C_0$ is a curve and $\{T_1 = 0\}$
cuts $C_0$ in finitely many points, so on a small enough $U$ the cut
locus is a finite set of points, and closure of a finite-type analytic
family of points is the union of the limits. It remains to show
$\overline{V_0} \cap U \cap \{T_1 = 0\} \subseteq \overline{C_0}$.

For the inclusion we proceed in two steps: we first replace the face
$\{T_1 = 0\}$ by the ambient closure $\overline{\Delta_\Gamma}$
(no analyticity beyond complex-analytic set theory is needed), and
then exclude extra one-dimensional components inside that closure by
a dimension count on $\Omega V_0$.

\emph{Step 1: closure of the boundary piece commutes with the stratum
decomposition.} Shrink $U$ so that
$\overline{\Delta_\Gamma} \cap U = \Delta_\Gamma^\circ \cap U$, where
$\Delta_\Gamma^\circ \supset \Delta_\Gamma$ is the union of the
finitely many strata of $\overline{\M}_g$ whose closures contain
$\Delta_\Gamma$ (the graphs obtained from $\Gamma$ by contracting
further nodes). Since
$\Delta_\Gamma^\circ$ is Zariski-locally closed in $\overline{\M}_g$
(it is open in its own closure, as is any stratum union of a
normal-crossings divisor decomposition), the set
$W_0 := \overline{V_0}\cap\Delta_\Gamma^\circ$ is a
complex-analytic subset of a normal analytic open set (the ambient
space minus finitely many strata disjoint from $\Delta_\Gamma$);
hence $\overline{W_0}\cap U$ is complex-analytic of pure dimension
$2$, and $W_0$ is Zariski-dense in $\overline{W_0}\cap U$. Every
point $P \in \overline{V_0}\cap\overline{\Delta_\Gamma}\cap U =
\overline{W_0} \cap U$ is therefore a limit of points of $W_0$. The
curve
\[
C_0 \;:=\; \overline{V_0}\cap\Delta_\Gamma \;=\; W_0 \cap
\Delta_\Gamma
\]
is Zariski-dense in $\overline{W_0}\cap U$: otherwise some
irreducible component $S$ of $\overline{W_0}\cap U$ would not meet
$\Delta_\Gamma$, but $S = \overline{S\cap W_0}$ by analyticity and
$W_0 \setminus \Delta_\Gamma = \varnothing$. Consequently every point
$P$ is a limit of points of $C_0$:
\[
\overline{V_0}\cap\overline{\Delta_\Gamma}\cap U
\;=\; \overline{\,V_0\cap\Delta_\Gamma\,}\cap U
\;=\; \overline{C_0}\cap U ,
\]
which is the load-bearing closure identity (the middle equality of
the lemma). Since $\Delta_{\Gamma\cup\mathcal N}\cap U \subset
\overline{\Delta_\Gamma}\cap U$ (on $U$ the strata adjacent to
$\Delta_\Gamma$ are nested in its closure), it follows that
\[
\overline{V_0}\cap U \cap \{T_1 = 0\}
\;=\; \overline{V_0}\cap\Delta_{\Gamma\cup\mathcal N}\cap U
\;\subseteq\; \overline{V_0}\cap\overline{\Delta_\Gamma}\cap U
\;=\; \overline{C_0}\cap U .
\]

\emph{Step 2 (only a consistency check; not needed for the
conclusion): the face $\mathcal F := \overline{V_0}\cap U\cap
\{T_1 = 0\}$ can in fact have no $1$-dimensional component.}
Indeed, suppose it had one, $\mathcal F_1$. Step 1 gives
$\mathcal F_1 \subset \overline{C_0}$, a curve, so $\mathcal F_1$
would be a boundary arc of $\overline{V_0}$ across the extra node
accumulated to $\overline{C_0}$. Pulling back the
$\GL_2^+(\R)$-action on $\Omega V_0$ through the face projection
$\Phi$, $\mathcal F_1$ supports a $1$-parameter family of stable
differentials in the boundary of $\Omega V_0$ whose
$\GL_2^+(\R)$-closure $\widehat{\mathcal F}_1$ is an affine
invariant manifold (Theorem~\ref{thm:EMM}(1)) projecting to the
closed invariant algebraic set
$\Phi(\widehat{\mathcal F}_1) \supset \mathcal F_1$
(Theorem~\ref{thm:EMM}(3)). Now $\Omega V_0$ has complex
dimension $3$ with $\dim_{\C}\Phi(\Omega V_0) = 2$
(Theorem~\ref{thm:Mukamel}(1)), so the generic fibre of
$\Phi|_{\Omega V_0}$ --- which contains the $\C^*$-orbit and, over
an interior point, equals the complex $1$-dimensional space of
disc-generating differentials --- has complex dimension $1$. Passing
to the boundary face over $\mathcal F_1$, the fibre can only be
smaller: a stable curve on a boundary divisor has finite
automorphism group, hence the stable differentials over a fixed
point of $\mathcal F_1$ in the boundary of $\Omega V_0$ form a
discrete union of $\C^*$-orbits, and
\[
\dim_{\C}\bigl(\Phi|_{\widehat{\mathcal F}_1}\bigr)^{-1}(x)
\;\le\; 1 \qquad (x \in \mathcal F_1 \ \text{generic}) .
\]
As
$\dim_{\C}\widehat{\mathcal F}_1 \le \dim_{\C}\Omega V_0 - 1 = 2$
(it lies in the boundary of the $3$-fold), this forces
$\dim_{\C}\Phi(\widehat{\mathcal F}_1) \le 1$: the projection of the
invariant closure is a curve, so $\mathcal F_1$ --- being a
$1$-dimensional irreducible component of that curve --- would have
to \emph{be} $\Phi(\widehat{\mathcal F}_1)$ itself, an invariant
curve. But Theorem~\ref{thm:DMMD}(1)-(4) identifies the boundary of
$\Omega V_0$ over the face through $\overline{C_0}$ with the central
Bers set of the boundary piece: its projection lies over the face of
$C_0$, i.e.\ over the discrete set $\overline{C_0}\cap\{T_1 = 0\}$
(stable-model uniqueness: the smoothing direction tangent to
$\overline{V_0}$ at a face point regenerates along $C_0$). A
$1$-dimensional $\Phi(\widehat{\mathcal F}_1)$ cannot lie over a
discrete set, contradiction. Hence $\mathcal F$ is
$0$-dimensional, and Step 1's inclusion is an equality with a
finite set. (This strengthens the inclusion but is not used in
Section~\ref{sec:conclusion}, which only needs
$\overline{V_0}\cap U\cap\{T_1=0\} \subseteq \overline{C_0}$.)
\end{proof}

\begin{prop}[Contradiction]\label{prop:contradiction}
Under the standing assumptions, the existence of such an $h$ is
impossible. Consequently the non-constant alternative cannot
occur, and $t \mapsto Y_t$ is constant for $|t| \ll 1$.
\end{prop}

\begin{proof}
By property (4) of Proposition~\ref{prop:regeneration}, the plumbing
coordinate $T_1$ of the collapsed node $\mathcal N$ vanishes identically
on $h(\{0\}\times\D)$ but is nonzero on $h(\varepsilon,s)$ for
$\varepsilon \neq 0$; hence the holomorphic map
$h_0 := h(0,\cdot) : \D \to \overline{\M}_g$ satisfies:
\begin{enumerate}[label=(\alph*)]
\item $h_0(s) \in \{T_1 = 0\} \cap \overline{V_0}$ for all $s$: indeed
$h_0(s) = \lim_{\varepsilon \to 0} h(\varepsilon, s)$ with
$h(\varepsilon,s) \in V_0$ (closedness of limits), and $T_1(h_0(s)) =
\lim T_1(h(\varepsilon,s)) = 0$;
\item $h_0$ is non-constant, with
$\partial_s (\pi_i \circ h_0)(0) = \dot Y_{t_*} \neq 0$.
\end{enumerate}
Let $U$ be the plumbing neighbourhood of $h_0(0)$ on which the
coordinates $(T_1,\dots,T_N; m)$ of Section~\ref{sec:tools} are defined;
after shrinking $\D$, $h_0(\D) \subset U$. Hypothesis (ii) applied with
$G = \Gamma$, $G' = \Gamma \cup \mathcal N$ (reading
$\Delta_{G'} \cap U = \Delta_\Gamma \cap \{T_1 = 0\} \cap U$ in the
plumbing model: the intersections living on $U$ of the two adjacent
strata differ exactly by $T_1 = 0$) gives
\[
\overline{V_0} \cap U \cap \{T_1 = 0\}
\;=\;
\overline{\,\overline{V_0} \cap \Delta_\Gamma\,} \cap \{T_1 = 0\} \cap U
\;=\;
\overline{C_0} \cap \{T_1 = 0\} \cap U
\;\subseteq\; \overline{C_0} \cap U.
\]
This identification at $t = 0$ is the content of
Lemma~\ref{lem:premise}: the first equality $\overline{\overline{V_0}
\cap \Delta_\Gamma} \cap \{T_1 = 0\} \cap U = \overline{V_0} \cap
\{T_1 = 0\} \cap U$ is a general analytic fact about closure commuting
with a proper divisor cut, and the inclusion $\subseteq
\overline{C_0}$ is proved there from the totally geodesic structure of
$V_0$ via Theorem~\ref{thm:EMM}(3) and Theorem~\ref{thm:DMMD}(1). Therefore
\[
h_0(\D) \;\subset\; \overline{V_0} \cap \{T_1 = 0\} \cap U
\;\subseteq\; \overline{C_0} \cap U .
\]
Now $\overline{C_0} \subset \Delta_\Gamma$ is a curve in the product
$\prod_j \M_{g_j,n_j}$ with $\pi_i(C_0) = \{Y_0\}$; the closure
$\overline{C_0}$ still projects under $\hat\pi_i$ to $\{Y_0\}$
(Lemma~\ref{lem:constclosure} below, noting that inside $U$ the
$i$-th component is unbroken, so the projection extends continuously as
$\hat\pi_i$ to $\overline{C_0}\cap U$). Hence $\hat\pi_i \circ h_0$
is a continuous, holomorphic map to the single point $Y_{t_*} =
\hat\pi_i(h_0(0))$, i.e.\ constant; but its derivative at $s=0$ is
$\partial_s(\hat\pi_i \circ h_0)(0) = \partial_s(\pi_i h)(0,0) \neq 0$,
a contradiction. \qedhere
\end{proof}

\begin{lem}[Constant projection on the closure]\label{lem:constclosure}
Let $C \subset \Delta_\Gamma$ be a closed analytic curve with
$\pi_i(C) = \{Y_0\}$, and let $\overline{C}$ be its closure in
$\overline{\M}_g$. Then $\pi_i(\overline{C} \cap \Delta_\Gamma)
= \{Y_0\}$; more precisely, if
$\hat \pi_i$ denotes the extension of $\pi_i$ to a neighbourhood of
$\overline{C}$ in which the $i$-th component of the stable curves remains
a single component (a plumbing neighbourhood), then
$\hat\pi_i(\overline{C}) = \{Y_0\}$.
\end{lem}

\begin{proof}
$\pi_i$ is the projection of a product, hence continuous and closed in a
neighbourhood of $\overline{C} \cap \Delta_\Gamma$; a continuous map
constant on a set is constant on its closure in the domain.
\end{proof}

\subsection{Globalization}

The local argument of the preceding sections (Sections
\ref{sec:slopes} through the contradiction
Proposition~\ref{prop:contradiction}) shows: after possibly shrinking
$\Delta$ around $0$, the holomorphic map $t \mapsto Y_t$ of
Lemma~\ref{lem:analytic} cannot be non-constant; hence there is a
neighbourhood $\Delta' \ni 0$ with $Y_t = Y_0$ for all $t \in \Delta'$.
(The argument is anchored at $t = 0$ because it uses the totally
geodesic structure of $V_0$; this is essential and corrects a previous
version which instead invoked a false re-centring of the whole argument
at arbitrary $t_0$ --- $V_{t_0}$ is not totally geodesic for $t_0 \ne
0$.)
Now $t \mapsto Y_t$ is a holomorphic map from the connected disc
$\Delta$ to $\M_{g_i,n_i}$ (Lemma~\ref{lem:analytic}), and it agrees
with the constant map $t \equiv Y_0$ on the open nonempty set
$\Delta'$. By the identity theorem for holomorphic maps between complex
manifolds (target Hausdorff), $Y_t = Y_0$ for all $t \in \Delta$. This
proves Theorem~\ref{thm:main}. \qed

\subsection{Where each hypothesis is used}
\begin{itemize}[leftmargin=1.8em]
\item $V_0$ totally geodesic and $2$-dimensional:
Theorem~\ref{thm:Mukamel} applies to $V_0$, making $\Omega V_0$ an
affine-invariant $3$-fold (this enters the dimension dichotomy of
Lemma~\ref{lem:membership} through $(O_0,q_0) \in \partial\Omega V_0$), and
gives the no-poles/cylinder-only degeneration of
Theorem~\ref{thm:Mukamel}(3) feeding Theorem~\ref{thm:DMMD}(4).
\item $\dim \partial V_t = 1$: makes $\partial V_t$ a curve so that the
plumbing slope analysis of Lemma~\ref{lem:slope} applies and so that the
degeneration data are well-defined stable differentials
(the limit along a $1$-dimensional boundary curve).
\item Hypothesis (i): smooth triviality gives the limiting
$\lim_{t\to0} y_t$, the analytic total space of Lemma~\ref{lem:analytic},
and the closed-analytic incidence level sets used in
Lemma~\ref{lem:membership} and Proposition~\ref{prop:constresidue}.
\item Hypothesis (ii): compatibility in deeper strata is used crucially in
Lemma~\ref{lem:slope} (the collapsed-node data and slopes at $t \ne 0$
equal those at $t=0$, computed in the deeper stratum for $V_0$) and in the
contradiction step (Proposition~\ref{prop:contradiction}, where
$\overline{V_0}\cap U\cap\{T_1=0\} \subseteq \overline{C_0}$ is exactly
hypothesis (ii)) and the globalization step.
\item $\pi_i(\partial V_t) = \{\mathrm{pt}\}$: produces the single-valued
map $t \mapsto Y_t$, and at $t=0$ gives
$\pi_i(\overline{C_0}) = \{Y_0\}$ used in
Proposition~\ref{prop:contradiction} via Lemma~\ref{lem:constclosure}.
\end{itemize}

\appendix

\section{Low-complexity factor cases}\label{sec:lowcomplexity}

The main argument (Sections \ref{sec:slopes}--\ref{sec:conclusion})
used the dynamics of $\PSL_2(\R)$ on strata and the exclusion of closed
totally geodesic surfaces (Theorem~\ref{thm:Mukamel}) in the factor
$\M_{g_i,n_i}$. Here we record the remaining low-complexity factors,
where the same conclusion holds either vacuously or by a mild
adaptation. There is nothing to change in
Lemma~\ref{lem:slope}, Lemma~\ref{lem:analytic},
Proposition~\ref{prop:constresidue} or the globalization: they use only
the plumbing coordinates, the hypotheses (i)--(ii), and Rado/Bishop ---
never the $\PSL_2(\R)$-structure.

\begin{description}[style=nextline]
\item[Factor a point: $(g_i,n_i) = (0,3)$.]
$\M_{g_i,n_i}$ is a single point. Then $\pi_i(\partial V_t)$ is
automatically the single point; $Y_t = Y_0$ and there is nothing to
prove.

\item[One-dimensional factor: $\dim_{\C}\M_{g_i,n_i} = 1$
(i.e.\ $(g_i,n_i) = (0,4)$, or $(1,1)$ as an orbifold: $\M_{1,1}$ is
a $1$-dimensional orbifold curve, not a point).]
Let $\{T_1 = 0\}$ be a boundary divisor meeting $C_0$
(Lemma~\ref{lem:slope}); by hypothesis (ii) (its $t = 0$ premise being
Lemma~\ref{lem:premise}),
$\overline{V_t} \cap \{T_1 = 0\} \cap U = \overline{C_t} \cap U$ for
all $t$, and this cut is a $0$-dimensional analytic cycle
$Z_t$ (a proper intersection of a surface with a divisor), varying
holomorphically in $t$ by hypothesis (i). Hence
$\pi_i(Z_t) \subset \M_{g_i,n_i}$ is a holomorphically varying
$0$-cycle; on the other hand $Z_t \subset \overline{C_t}$, and
$\pi_i(\overline{C_t})$ is the singleton $\{Y_t\}$
(Lemma~\ref{lem:constclosure}), so $\pi_i(Z_t) = \{Y_t\}^{\oplus \deg
Z_t}$ as a $0$-cycle concentrated at $Y_t$. The support
$t \mapsto Y_t$ of a holomorphic family of $0$-cycles of constant
degree concentrated at a point is a holomorphic map (it is the support
of the image cycle); by Lemma~\ref{lem:analytic} it is the map already
denoted $t \mapsto Y_t$. Since $Z_t$ is a flat deformation of
$Z_0$, which is supported at $Y_0$, and a flat family of $0$-cycles
of constant degree with singleton support has constant support (the
cycle map is continuous and the support of $Z_0$ is $Y_0$), we get
$Y_t = Y_0$ for all small $t$, hence for all $t$ by the identity
theorem. (This argument does not use any
$\PSL_2(\R)$-structure, so it applies to every factor of dimension
$1$, including $(0,4)$ and $(1,1)$.)

\item[Genus zero with $n_i \ge 5$: $(g_i,n_i) = (0,n_i)$, $n_i \ge 5$.]
The argument proceeds verbatim as in the main case. The strata of
meromorphic quadratic differentials on $\mathbb P^1$ still carry the
complex-affine period-coordinate structure, the
$\GL_2^+(\R)$-action, and the plumbing/regeneration theory of
Section~\ref{sec:tools} (the references cited there remain valid in
genus zero), and the membership, regeneration, and contradiction
arguments of Sections~\ref{sec:slopes}--\ref{sec:conclusion} apply
verbatim. The exclusion of a $2$-dimensional $\Phi$-image is not
needed here: Proposition~\ref{prop:contradiction} is
$g_i$-independent, using only plumbing coordinates and
hypothesis (ii). Hence $Y_t = Y_0$.

\item[Genus two: $g_i = 2$, $n_i \ge 0$.]
$\M_2$ contains McMullen's Hilbert modular surfaces
\cite{McMullenHMS}: genuine totally geodesic surfaces, projections of
affine invariant $3$-folds with $2$-dimensional base. If
$\Phi(B_0)$ were such a surface $W_H \subset \M_{2,n_i}$, then the
boundary face of $\Omega V_0$ over the corresponding face of
$\partial \M_g$ would dominate $W_H$; but $\Omega V_0$ fibres over
$V_0 \subset \M_g$ with $1$-dimensional fibres away from the boundary,
and taking closures identifies the fibre of $\Phi|_{B_0}$ generically
with the disc bundle of the boundary stratum: a
$2$-dimensional $\Phi(B_0)$ is possible, and in that case
$\{(O_t,q_t)\}$ varies in a closed totally geodesic surface. (Here
Input~(Mem) is still assumed, as in Section~\ref{sec:slopes}.)
However, hypothesis (ii) and the regeneration argument of
Section~\ref{sec:conclusion} apply literally: the conclusions of
Proposition~\ref{prop:contradiction} (which uses only plumbing
coordinates and hypothesis (ii), not the value of $g_i$) force $Y_t =
Y_0$ in this case as well. (The only place in Sections
\ref{sec:slopes}--\ref{sec:conclusion} where $g_i$ enters is through
Theorem~\ref{thm:Mukamel}(2), used to exclude a priori the possibility
that $\Phi(B_0)$ is a totally geodesic surface; replacing this
exclusion requires no further argument, since the contradiction of
Proposition~\ref{prop:contradiction} is deduced independently of the
nature of $\Phi(B_0)$.)
\end{description}

\section{Status of the main statement and the single external input}
\label{sec:status}

For the convenience of the reader we collect the logical structure.
Theorem~\ref{thm:main} is proved from: (i) the hypotheses of the problem
statement; (ii) the four quoted inputs
(Theorems~\ref{thm:EMM}, \ref{thm:Mukamel}, \ref{thm:DMMD} and
Proposition~\ref{prop:regeneration}); and (iii) the single
\emph{additional analytic input}, Input~(Mem)
(p.\ \pageref{input:membership}), whose content is: the limiting stable
differentials of the whole deformation remain in the affine manifold
generated from the $V_0$ boundary piece. Input~(Mem) is expected to
follow from the full strength of the boundary theory of affine invariant
manifolds applied to the total space of the family $\mathcal V$;
isolating it here makes the logical dependencies explicit and confines
the non-elementary analytic content of the argument to a single,
precisely stated place.

\emph{Consistency certificate with the problem statement.} The
conclusion of Theorem~\ref{thm:main} is exactly the conclusion asked
for in the problem: $\pi_i(\partial V_t) = \{Y_0\}$ literally for
\emph{all} $t \in \Delta$, not merely for small $t$ (the identity
theorem of Section~\ref{sec:conclusion} lifts the local constancy near
$t = 0$ to the whole connected disc). The hypotheses of
Theorem~\ref{thm:main} match the problem statement term-for-term; the
only logical residue relative to the problem's literal request is the
single unproved analytic input Input~(Mem), isolated and labelled ---
not concealed.

\begin{rmk}[Self-diagnosis: where the missing anchor sits]
At the level of stable curves, the missing anchor of the argument is
exactly the $t \neq 0$ form of the premise of hypothesis (ii) ---
Lemma~\ref{lem:premise} supplies that premise only at $t = 0$, which is
why the contradiction of Proposition~\ref{prop:contradiction} is staged
at $t = 0$ via regeneration into $V_0$. Input~(Mem) sits precisely one
$\Phi$-level above this: it is the same compatibility, lifted from
stable curves to the quadratic differentials and their affine-linear
invariance equations.
\end{rmk}

\begin{thebibliography}{9}

\bibitem{EM} A.~Eskin and M.~Mirzakhani,
\emph{Invariant and stationary measures for the
$\mathrm{SL}(2,\mathbb{R})$ action on moduli space},
Publ.\ Math.\ Inst.\ Hautes \'Etudes Sci.\ \textbf{127} (2018), 95--324.

\bibitem{EMM} A.~Eskin, M.~Mirzakhani, and A.~Mohammadi,
\emph{Isolation, equidistribution, and orbit closures for the
$\mathrm{SL}(2,\mathbb{R})$ action on moduli space},
Ann.\ of Math.\ (2) \textbf{182} (2015), no.~2, 673--721.

\bibitem{Filip} S.~Filip,
\emph{Splitting mixed Hodge structures over affine invariant manifolds},
Ann.\ of Math.\ (2) \textbf{183} (2016), no.~2, 681--713; and
\emph{Semisimplicity and rigidity of the Kontsevich--Zorich cocycle},
Invent.\ Math.\ \textbf{205} (2016), no.~3, 617--670.

\bibitem{Mukamel} R.~Mukamel,
\emph{Totally geodesic subvarieties of the moduli space of Riemann
surfaces}; cited in the text in the restated form of
Theorem~\ref{thm:Mukamel}(1)--(2). \emph{Bibliographic note}: the
precise journal reference could not be verified without external lookup
in this environment; the statement used is quoted self-containedly in
the text.

\bibitem{McMullenHMS} C.~McMullen,
\emph{Dynamics of $\mathrm{SL}_2(\mathbb{R})$ over moduli space in genus
two}, Ann.\ of Math.\ (2) \textbf{165} (2007), no.~2, 397--456.

\bibitem{SW} J.~Smillie and B.~Weiss,
\emph{Characterizations of lattice surfaces}, Invent.\ Math.\
\textbf{180} (2010), no.~3, 535--557.

\bibitem{MW} N.~Mazor and A.~Wright,
\emph{The boundary of an affine invariant submanifold},
preprint (2024).

\bibitem{BCGGM} M.~Bainbridge, D.~Chen, Q.~Gendron, S.~Grushevsky, and
M.~M\"oller,
\emph{Compactification of strata of abelian differentials},
Duke Math.\ J.\ \textbf{167} (2018), no.~12, 2347--2416.

\end{thebibliography}

\end{document}
